%%%PAGE 296 rot=0 legibility=good summary=自感(通电/断电自感电路、理想线圈)；多用电表欧姆挡测R/C/L/D；斜轨导体棒能量守恒与电量、热量分配
%%%BLOCK id=296-01 ch=12 sec=自感 kind=knowledge alt=none cont=none y=0.03-0.31
\textbf{自感}\quad \unclear{感}阻碍电流变化\quad $i$增反减同

% 左图线圈旁另有小字标注（字迹不清，形如“R_L+r→0”）
%FIG p296_a 0.085 0.068 0.335 0.19
\notefig[0.4]{figures/p296_a.png}

%FIG p296_b 0.455 0.068 0.71 0.19
\notefig[0.4]{figures/p296_b.png}

\textbf{通电自感}\quad 闭A后亮\quad B立即亮

断电自感\quad AB过一段才熄灭

$I_B>I_A$\quad A闪后逐渐熄灭
%%%END
%%%BLOCK id=296-02 ch=12 sec=自感 kind=knowledge alt=none cont=none y=0.31-0.44
\textbf{理想线圈}
\begin{itemize}
\item 通电时相当于电断路 % CHECK: “电断路”，“电”字疑为多写，原样转录
\item 断电时相当于电源
\item 稳定时相当于导线
\end{itemize}
%%%END
%%%BLOCK id=296-03 ch=10 sec=多用电表·欧姆挡测元件 kind=knowledge exp=1 alt=none cont=none y=0.43-0.58
黑接正极\quad 正向电阻

\begin{tabular}{@{}l@{\hspace{2.5em}}l}
R & 不变 \\
C & 先向右后向左 \\
L & 缓慢向右偏 \\
D & 正负极电阻相差大
\end{tabular}
%%%END
%%%BLOCK id=296-04 ch=12 sec=电磁感应·能量与热量分配 kind=problem alt=06 cont=none y=0.60-0.90
\begin{problem}
% 图中标注（已在图里）：R_1、2r、R_2、r（棒）、37°、x_0
%FIG p296_c 0.075 0.63 0.357 0.81
\notefig[0.35]{figures/p296_c.png}

\begin{solution}
% 原稿：$W_f$ 似被划去，其正上方写有 $-\mu mgx_0\cos\alpha$；$W_f$ 与 $mgx_0\sin\alpha$ 之间是一个“·”
\[
\overset{-\mu mgx_0\cos\alpha}{W_f}\cdot mgx_0\sin\alpha=\frac{1}{2}mv_{\max}^2+Q
\] % CHECK: “W_f”与“mgx_0 sinα”之间为“·”，物理上应为“+”

\textbf{电量、热量分配}

求 $Q_{\text{总}}\Rightarrow Q_{\text{外}}$\ \ $\left(\dfrac{Q_{\text{外}}}{Q_{\text{总}}}=\dfrac{R}{R+r}\right)$\ \ $\Rightarrow Q\propto Q_1\ Q_2$ % CHECK: “Q ∝ Q_1 Q_2”中间符号不清

（括号上方小字：与 $R$ 成正比。右侧并列两行小字：与 $R$ 成反比（并联）；与 $R$ 成正比（串联））

并联 $Q$ 与 $R$ 成反比\quad $\dfrac{U^2}{R}$

串联 $Q$ 与 $R$ 成正比

\[
q=\frac{BLx}{R_{\text{总}}}
\]

$q_1$、$q_2$ 与电阻成反比。
\end{solution}
\end{problem}
%%%END
%%%PAGEEND 296
